% Part:sets-functions-relations
% Chapter: size-of-sets
% Section: comparing-sizes

\documentclass[../../../include/open-logic-section]{subfiles}

\begin{document}

\olfileid{sfr}{siz}{car}

\olsection{వేర్వేరు పరిమాణాల సమితులు, కాంటర్ ప్రమేయం}

\begin{explain}
రెండు సమితులు ఒకే పరిమాణంలో ఉన్నాయనే భావనకు కచ్చితమైన ప్రకటనను
ఇచ్చాం. ఒక సమితి మరొకదాని కంటే చిన్నదనే భావనకూ కచ్చితమైన ప్రకటనను
ఇవ్వవచ్చు. ``కంటే చిన్నది (లేదా సమసంఖ్యాకం)'' అనే మన నిర్వచనానికి
సమితుల మధ్య \tetoken{ఒక ద్విగుణ ప్రమేయానికి}{!!a{bijection}} బదులు, మొదటి సమితి నుంచి రెండవదానికి
\tetoken{ఒక ఒకటి-ఒకటి ప్రమేయం}{!!a{injection}} అవసరం. అటువంటి ప్రమేయం ఉంటే, మొదటి సమితి పరిమాణం
రెండవ సమితి పరిమాణానికి తక్కువ లేదా సమానం. సహజంగా చూస్తే, ఒక సమితి
నుంచి మరొకదానికి \tetoken{ఒక ఒకటి-ఒకటి ప్రమేయం}{!!a{injection}} ఉండటం వల్ల ఆ ప్రమేయపు పరిధిలో దాని
ప్రవేశంలో ఉన్నన్ని \tetoken{మూలకాలు}{!!{element}s} కనీసం ఉంటాయని హామీ వస్తుంది; ఎందుకంటే
ప్రవేశంలోని ఏ రెండు \tetoken{మూలకాలు}{!!{element}s} కూడా పరిధిలోని ఒకే \tetoken{మూలకానికి}{!!{element}}
పంపబడవు.
\end{explain}

\begin{defn}
\tetoken{ఒక ఒకటి-ఒకటి ప్రమేయం}{!!a{injection}} $f \colon A \to B$ ఉంటేనే $A$, $B$ను
\emph{పరిమాణంలో మించదు} అంటాం; దీన్ని $\cardle{A}{B}$ అని రాస్తాం.
\end{defn}

ఇది స్వావర్తన, సంక్రామక సంబంధమనీ, కానీ సౌష్ఠవ సంబంధం కాదనీ స్పష్టం
(దీన్ని అభ్యాసంగా వదిలాం). ఒక సమితి మరొకదాని కంటే (కచ్చితంగా)
చిన్నదని తెలిపే భావనను కూడా పరిచయం చేయవచ్చు.

\begin{defn}
\tetoken{ఒక ఒకటి-ఒకటి ప్రమేయం}{!!a{injection}}~$f\colon A \to B$ ఉండి, \tetoken{ద్విగుణ ప్రమేయం}{!!{bijection}}~$g\colon A
\to B$ ఏదీ లేకపోతేనే $A$, $B$ కంటే \emph{పరిమాణంలో చిన్నది} అంటాం;
దీన్ని $\cardless{A}{B}$ అని రాస్తాం. అంటే, $\cardle{A}{B}$ మరియు
$\cardneq{A}{B}$.
\end{defn}

ఈ సంబంధం అస్వావర్తన, సంక్రామకం అని స్పష్టం. (దీన్ని అభ్యాసంగా
వదిలాం.) ఈ సంకేతంతో, $A$ \tetoken{లెక్కించదగిన}{!!{enumerable}} కావడానికి అవసరమైన మరియు
సరిపడే షరతు $\cardle{A}{\Nat}$ అవడమనీ, $A$ \tetoken{లెక్కించలేని}{!!{nonenumerable}}
కావడానికి అవసరమైన మరియు సరిపడే షరతు $\cardless{\Nat}{A}$ అవడమనీ
చెప్పవచ్చు. దీనివల్ల
\oliflabeldef{sfr:siz:nen-alt:thm:nonenum-pownat}{%
\olref[sfr][siz][nen-alt]{thm:nonenum-pownat}ను
$\cardless{\Nat}{\Pow{\Nat}}$ అనే గమనింపుగా}{%
\olref[sfr][siz][nen]{thm:nonenum-pownat}ను
$\cardless{\PosInt}{\Pow{\PosInt}}$ అనే గమనింపుగా} తిరిగి చెప్పవచ్చు.
నిజానికి, ఈ చివరి విషయం \emph{పూర్తిగా సార్వత్రికం} అని
\citet{Cantor1892} నిరూపించాడు:

\begin{thm}[కాంటర్]\ollabel{thm:cantor}
ఏ సమితి $A$కైనా $\cardless{A}{\Pow{A}}$.
\end{thm}

\begin{proof}
$f(x) = \{x\}$ అనే చిత్రణ $f \colon A \to \Pow{A}$ ఒక
\tetoken{ఒకటి-ఒకటి ప్రమేయం}{!!a{injection}}. ఎందుకంటే $x \neq y$ అయితే సమితుల సమానత్వ సూత్రం వల్ల
$\{x\} \neq \{y\}$ కూడా; కాబట్టి $f(x) \neq f(y)$. అందువల్ల
$\cardle{A}{\Pow{A}}$.

\begin{editorial}
\olref[nen]{sec} ఉంటే వివరణాత్మక నిరూపణను, లేకపోతే
\olref[nen-alt]{sec}కు సరిపోలే వేగవంతమైన నిరూపణను ఇస్తాం.
\end{editorial}

\oliflabeldef{sfr:siz:nen:sec}{%
  ఇప్పుడు \tetoken{ఒక సంగ్రస్త}{!!a{surjective}} ప్రమేయం~$g\colon A \to \Pow{A}$ ఉండజాలదని,
  ఇక \tetoken{ఒక ద్విగుణ}{!!a{bijective}} ప్రమేయం సంగతి చెప్పనక్కరలేదని చూపి,
  $\cardneq{A}{\Pow{A}}$ అని తేలుస్తాం. $g\colon A \to \Pow{A}$
  అనుకుందాం. $g$ సంపూర్ణం కనుక, ప్రతి $x \in A$ ఏదో ఉపసమితి
  $g(x) \subseteq A$కు పంపబడుతుంది. $g$ సంగ్రస్తం కాలేదని చూపవచ్చు.
  అందుకోసం, నిర్వచనం వల్లనే $g$ పరిధిలో ఉండజాలని ఒక ఉపసమితి
  $\overline{A} \subseteq A$ను నిర్వచిద్దాం. ఇలా తీసుకోండి:
  \[
  \overline{A} = \Setabs{x \in A}{x \notin g(x)}.
  \]
  ప్రతి $x \in A$కూ $g(x)$ నిర్వచించబడింది కనుక, $\overline{A}$
  స్పష్టంగా $A$ యొక్క సునిర్వచిత ఉపసమితి. కానీ అది $g$ పరిధిలో
  ఉండజాలదు. ఏదైనా $x \in A$ తీసుకుందాం; $\overline{A} \neq g(x)$
  అని చూపుతాం. $x \in g(x)$ అయితే, అది $x \notin g(x)$ను
  సంతృప్తిపరచదు; కాబట్టి $\overline{A}$ నిర్వచనం వల్ల
  $x \notin \overline{A}$. $x \in \overline{A}$ అయితే,
  $\overline{A}$ నిర్వచక ధర్మాన్ని అది సంతృప్తిపరచాలి; అంటే,
  $x \in A$, $x \notin g(x)$. $x$ యాదృచ్ఛికం కనుక, ప్రతి
  $x \in A$కు $x \in g(x)$ అవడానికి అవసరమైన మరియు సరిపడే షరతు
  $x \notin \overline{A}$ అవడం; అందువల్ల
  $g(x) \neq \overline{A}$. మరోలా చెప్పాలంటే, $\overline{A}$
  $g$ పరిధిలో ఉండజాలదు; కాబట్టి $g$ సంగ్రస్తం కాదు.
  \sourcecorrection{OLSIZ-007}{యాదృచ్ఛికంగా తీసుకున్న $x$ గురించి
  తేల్చే చోట ఆంగ్ల మూలం పొరపాటుగా ప్రతి $x\in\overline{A}$కు అని
  పరిమితం చేస్తుంది. వికర్ణ వాదనకు అవసరమైనట్లుగా, ఇక్కడ ప్రతి
  $x\in A$కు అని సరిచేశాం.}}{ఇంకా $\cardneq{A}{\Pow{A}}$ అని
  చూపాలి. విరోధం కోసం $\cardeq{A}{\Pow{A}}$ అనుకుందాం; అంటే, ఏదో
  \tetoken{ద్విగుణ ప్రమేయం}{!!{bijection}} $g \colon A \to \Pow{A}$ ఉంది. ఇప్పుడు దీన్ని
  పరిశీలించండి:
  \[
    D = \Setabs{x \in A}{x \notin g(x)}
  \]
  $D \subseteq A$ కనుక $D \in \Pow{A}$. $g$ \tetoken{ఒక ద్విగుణ ప్రమేయం}{!!a{bijection}} కనుక,
  $g(y) = D$ అయ్యే ఏదో $y \in A$ ఉంటుంది. కానీ ఇప్పుడు:
  \[
    y \in g(y) \text{ iff } y \in D \text{ iff } y \notin g(y).
  \]
  ఇది విరోధం; కాబట్టి $\cardneq{A}{\Pow{A}}$.}{}
\end{proof}

\begin{explain}
\oliflabeldef{sfr:siz:nen:thm:nonenum-pownat}{
  \olref{thm:cantor} నిరూపణను
  \olref[nen]{thm:nonenum-pownat} నిరూపణతో పోల్చడం బోధప్రదం. అక్కడ
  $\PosInt$ ఉపసమితుల ఏ జాబితా $Z_1$, $Z_2$, \dots{} ఇచ్చినా, ఆ
  జాబితాలో ఉండదని హామీ ఉన్న సంఖ్యల సమితి $\overline{Z}$ను
  నిర్మించవచ్చని చూపాం. ప్రతి $n \in \PosInt$కూ $n \in Z_n$
  అవడానికి అవసరమైన మరియు సరిపడే షరతు $n \notin \overline{Z}$ అవడం
  కనుక అది జాబితాలో ఉండదని హామీ వచ్చింది. ఈ విధంగా, $Z_n$,
  $\overline{Z}$లలో ఒకదానికి \tetoken{ఒక మూలకం}{!!a{element}} అయి, మరొకదానికి కాని సంఖ్య
  ఎప్పుడూ ఉంటుంది. ఇక్కడ కూడా అదే ఆలోచనను అనుసరిస్తాం; అయితే సూచీలు
  $n$ ఇప్పుడు $\PosInt$కు బదులు $A$ యొక్క \tetoken{మూలకాలు}{!!{element}s}. ప్రతి
  $x \in A$కూ $x \in g(x)$ అవడానికి అవసరమైన మరియు సరిపడే షరతు
  $x \notin \overline{A}$ అవడం వల్ల, $\overline{A}$ ప్రతి $g(x)$కు
  భిన్నంగా ఉండేలా నిర్వచించబడింది. మళ్లీ, $g(x)$, $\overline{A}$లలో
  ఒకదానికి \tetoken{ఒక మూలకం}{!!a{element}} అయి, మరొకదానికి కాని $A$ యొక్క
  \tetoken{ఒక మూలకం}{!!a{element}} ఎప్పుడూ ఉంటుంది. అందువల్ల $\overline{Z}$, $Z_1$,
  $Z_2$, \dots జాబితాలో ఉండలేనట్లే, $\overline{A}$ కూడా $g$ పరిధిలో
  ఉండజాలదు.}{}

\oliflabeldef{sfr:siz:nen-alt:thm:nonenum-pownat}{
  \olref{thm:cantor} నిరూపణను
  \olref[nen-alt]{thm:nonenum-pownat} నిరూపణతో పోల్చడం బోధప్రదం.
  అక్కడ $\Nat$ ఉపసమితుల ఏ జాబితా $N_0$, $N_1$, $N_2$, \dots{}
  ఇచ్చినా, ఆ జాబితాలో ఉండదని హామీ ఉన్న సంఖ్యల సమితి $D$ను
  నిర్మించవచ్చని చూపాం. ప్రతి $n \in \Nat$కూ $n \in N_n$
  అవడానికి అవసరమైన మరియు సరిపడే షరతు $n \notin D$ అవడం వల్ల అది
  జాబితాలో ఉండదని హామీ వచ్చింది. ఇక్కడా అదే ఆలోచనను అనుసరిస్తాం;
  అయితే సూచీలు $n$ ఇప్పుడు $\Nat$కు బదులు $A$ యొక్క \tetoken{మూలకాలు}{!!{element}s}.
  ప్రతి $x \in A$కూ $x \in g(x)$ అవడానికి అవసరమైన మరియు సరిపడే
  షరతు $x \notin D$ అవడం వల్ల, $D$ ప్రతి $g(x)$కు భిన్నంగా ఉండేలా
  నిర్వచించబడింది.}{}

ఈ నిరూపణను రస్సెల్ విరోధాభాస నిరూపణ
\olref[sfr][set][rus]{thm:russells-paradox}తో పోల్చడం కూడా విలువైనదే.
నిజానికి, కాంటర్ ప్రమేయమే రస్సెల్ సొంత విరోధాభాసానికి ప్రేరణ.
\end{explain}

\begin{prob}
  ఏ సమితి~$A$కైనా \tetoken{ఒక ఒకటి-ఒకటి ప్రమేయం}{!!a{injection}} $g\colon \Pow{A} \to A$ ఉండజాలదని
  చూపండి. సూచన: $g\colon \Pow{A} \to A$ \tetoken{ఒకటి-ఒకటి}{!!{injective}} అనుకుందాం.
  $D = \Setabs{g(B)}{B \subseteq A \text{ and } g(B) \notin B}$ను
  పరిశీలించండి. $x = g(D)$గా తీసుకోండి. $g$ \tetoken{ఒకటి-ఒకటి}{!!{injective}} అనే
  వాస్తవాన్ని వాడి ఒక విరోధాన్ని ఉత్పాదించండి.
\end{prob}

\end{document}
